\documentclass{kpi-lab}
\usepackage{biblatex}
\usepackage{enumitem}
\usepackage{longtable}
\usepackage{namedlisting}
\addbibresource{sources.bib}

\title{Лабораторна робота \textnumero 3\\
  Реалізація основних асиметричних криптосистем
}
\author{ФБ-51мн Подолянко Тимофій}
\date{}

\begin{document}
\maketitle
\setcounter{tocdepth}{1}
\tableofcontents
\bigskip

\textbf{Мета роботи} ---
дослідити принци побудови та застосування сучасних
асиметричних криптосистем,
отримати практичні навички використання механізмів
цифрового підпису та узгодження ключів,
проаналізувати особливості безпечного застосування
асиметричних криптографічних алгоритмів.

\chapter{Дослідження принципів асиметричної криптографії}
\label{chap:asymcrypt-principles}

\begin{quote}\small\it
  У чому полягає основний принцип асиметричної криптографії
  та які функції виконують відкритий і закритий ключі?
\end{quote}

Асиметрична криптографія спирається на алгоритми, які застосовують два
пов'язаних ключа: публічний та секретний. Властивістю пари
ключів є те, що знаходження секретного ключа, який
відповідає публічному ключу, є обчислювально складною
задачею \cite{nist-sp-800-57pt1r5}.

\begin{quote}\small\it
  Які принципові відмінності між асиметричним шифруванням,
  цифровим підписом та узгодженням спільного секрету (key agreement)?
\end{quote}

Асиметричне шифрування, цифровий підпис та узгодження
спільного секрету виконують різні задачі, хоча й гарантії
стійкості захисту кожного ґрунтується на тих самих
обчислювальних проблемах, і частково використовуються ті
самі криптопримітиви.

Асиметричне шифрування дозволяє гарантувати
конфіденційність інформації. Будь-хто може захистити
інформацію за допомогою публічного а розкриття захищеної інформації
можливе лише власником відповідного секретного ключа.

Цифровий підпис гарантує цілісність захищеної інформації, а також
автентифікацію джерела підпису (та, за достатніх, зокрема
правових, умов, неможливість відмови від авторства)
\cite{nist-sp-800-57pt1r5}.

Узгодження спільного секрету гарантує створення
секретного значення, яке є функцією від внеску обох учасників
і не може бути передбачене жодною стороною незалежно від
внеску іншої \cite{nist-sp-800-56ar3}.
Додатково гарантується, що жодна сторона не може
безпосередньо контролювати обчислене секретне значення.

\begin{quote}\small\it
  Які задачі вирішують RSA, ECDSA/EdDSA та ECDH/X25519?
  Чому ці алгоритми не можна вважати взаємозамінними?
\end{quote}

Примітиви RSA шифрування та дешифрування дозволяють,
відповідно, здійснювати шифрування в цій асиметричній
криптосистемі \cite{nist-sp-800-56ar3}.
ECDSA, EdDSA реалізують цифровий підпис.
ECDH, X25519 реалізують узгодження спільного секрету.
Це різні задачі, та коректна реалізація кожного алгоритма повинна
враховувати низку нюансів специфічних для задачі. Так,
наприклад, для цифрових підписів ECDSA секретне значення
для повідомлення не повинне повторюватися для двох різних повідомлень,
інакше тривіальним стає розкриття секретного ключа
\cite{cryptoeprint:2025/654, MECANISMES-CRYPTO-1.0}.

\begin{quote}\small\it
  Чому довжини ключів RSA та криптосистем на еліптичних
  кривих не можна безпосередньо порівнювати для оцінювання
  рівня криптографічної стійкості?
\end{quote}

Для коректної оцінки рівня криптографічної стійкості
необхідно враховувати математичні властивості примітивів та застосовні
обчислювальні алгоритми, апаратні засоби, а загалом на таку
оцінку впливає й інша інформація, відома нападнику.
Тому, наприклад, для ключів RSA з модулем довжиною 2048 бітів
рівень стійкості захисту оцінюється лише в 112 бітів,
а для ECDSA аналогічна оцінка досягається вже при 224-255
бітових ключах (відповідному розміру порядку базової точки).
\cite{nist-sp-800-57pt1r5}.

\begin{quote}\small\it
  Для чого при використанні RSA застосовуються спеціальні
  схеми доповнення (padding)? Які ризики створює
  неправильний вибір або використання padding?
\end{quote}

Padding усуває детермінізм результату RSA шифрування, що
унеможливлює, чи принаймні значно ускладнює,
атаки на основі знання шифротексту (наприклад, простий
частотний аналіз).

Небезпечне застосування доповнення може зробити систему вразливою до
варіантів атак з обраним шифротекстом. Як приклад можна навести
атаку з обраним шифротекстом на PKCS\#1 v1.5 RSA
\cite{Bleichenbacher1998}

Padding для цифрового підпису, такий як RSASSA-PSS,
зберігає детермінізм,
але включає в підписане повідомлення додаткову інформацію
про власне процес підпису, яка також повинна перевірятися
під час перевірки підпису.
\cite{fips186-5}.

\begin{quote}\small\it
  Яку роль виконує nonce при формуванні цифрового підпису
  ECDSA? Які наслідки може мати повторне використання
  одного nonce для різних повідомлень?
\end{quote}

Внесення nonce у криптосистему додає в систему рівнянь,
якою можливо було б визначити секретний ключ за множиною
пар  повідомлень і відповідних коректних підписів,
додаткову змінну.
Тоді стають можливими безліч розв'язків такої системи,
тобто відновити секретний ключ \enquote{в обхід} задачі
пошуку дискретного логарифма не вдасться.
Для цифрових підписів ECDSA nonce є секретним і не повинен
повторюватися для двох різних повідомлень
\cite{10.1007/3-540-39568-7_2}.

Повторне використання nonce робить тривіальним відновлення
секретного ключа з двох різних повідомлень, підписаних з
однаковим nonce \cite{cryptoeprint:2025/654}.

\begin{quote}\small\it
  Чому асиметричне шифрування зазвичай не використовується
  безпосередньо для захисту великих обсягів даних? Що таке
  гібридна криптосистема?
\end{quote}

Через особливості примітивів в основі асиметричного шифрування,
процедури шифрування й дешифрування дуже повільні у порівнянні з
симетричними алгоритмами. Крім цього, розмір блоку для
операції обмежений довжиною модуля \cite{nist-sp-800-56br2}.

Гібридні криптосистеми, за визначенням \cite{rfc4949}, поєднують
застосування симетричних та асиметричних алгоритмів
для отримання переваг обох підходів: \enquote{керування
  ключами як в асиметричних системах, а швидкість
(де)шифрування як в симетричних} \cite{rfc9180}.

\enquote{Гібридна криптосистема} також використовується для позначення
асиметричної криптосистеми, яка застосовує декілька
алгоритмів, серед яких є традиційні та постквантово-стійкі
алгоритми \cite{MECANISMES-CRYPTO-1.0}.

\begin{quote}\small\it
  Які вимоги висуваються до зберігання та використання
  закритого ключа асиметричної криптосистеми?
\end{quote}

Секретний ключ є критичним об'єктом асиметричних криптосистем.
Зберігання та використання секретних ключів регламентуються
у \cite{nist-sp-800-57pt1r5} згідно з визначеною моделлю
життєвого циклу ключа. Загалом заходи захисту включають
захищену генерацію, зберігання; контроль та аудит доступу і
застосування;
обмеження на цільове використання (в тому сенсі, що певний
  ключ не повинен застосовуватися для кількох різних
криптографічних задач \cite{MECANISMES-CRYPTO-1.0});
дотримання криптоперіоду;
безпечне і гарантоване знищення ключового  матеріалу після
закінчення використання.

\chapter{Практичне дослідження цифрового підпису}

Розглянемо програму на Zig v0.17.1, яка реалізує генерацію ключів,
формування і перевірку цифрового підпису текстових повідомлень.
У лістингу \ref{listing:2.1} наведено текст програми.
Криптографічний алгоритм реалізований в стандартній бібліотеці мови.

\namedfilelisting{zig}{2.1.zig}
{Програма на Zig для формування і перевірки цифрового підпису Ed25519}
\label{listing:2.1}

Генеруємо ключову пару:
\begin{minted}{shell}
./2.1.zig generate
\end{minted}
\begin{minted}{output}
Public key: 8fcf6caa38ec7e6906aa42add93c12a7eb541730d6316225f9ab010de72f00af
Private key: 70fa8f80dd52925148ae77635f5660ab4cd8fe462bc00f620042ab9eacdb560f8fcf6caa38ec7e6906aa42add93c12a7eb541730d6316225f9ab010de72f00af
\end{minted}

Підписуємо \enquote{тестове повідомлення}:
\begin{minted}{shell}
./2.1.zig sign \
  70fa8f80dd52925148ae77635f5660ab4cd8fe462bc00f620042ab9eacdb560f8fcf6caa38ec7e6906aa42add93c12a7eb541730d6316225f9ab010de72f00af \
  'тестове повідомлення'
\end{minted}
\begin{minted}{output}
Signature: 038a4328e0c79249159080f0444edd04a2b22bbdc2c9dbdf6465e201b9eac7faf22f6ab7c4bf94862d9c00851cc6f62da088bfcbba4265fb5297485db0ff9f01
\end{minted}

Перевіряємо підпис за допомогою відкритого ключа:
\begin{minted}{shell}
./2.1.zig verify \
  8fcf6caa38ec7e6906aa42add93c12a7eb541730d6316225f9ab010de72f00af \
  'тестове повідомлення' \
  038a4328e0c79249159080f0444edd04a2b22bbdc2c9dbdf6465e201b9eac7faf22f6ab7c4bf94862d9c00851cc6f62da088bfcbba4265fb5297485db0ff9f01
\end{minted}
\begin{minted}{output}
Signature valid!
\end{minted}

Змінимо повідомлення і спробуємо перевірити:
\begin{minted}{shell}
./2.1.zig verify \
  8fcf6caa38ec7e6906aa42add93c12a7eb541730d6316225f9ab010de72f00af \
  'тестове Повідомлення' \
  038a4328e0c79249159080f0444edd04a2b22bbdc2c9dbdf6465e201b9eac7faf22f6ab7c4bf94862d9c00851cc6f62da088bfcbba4265fb5297485db0ff9f01
\end{minted}
\begin{minted}{output}
INVALID SIGNATURE!
error: SignatureVerificationFailed
\end{minted}

Спробуємо перевірити підпис за допомогою іншого ключа:
\begin{minted}{shell}
./2.1.zig verify \
  c52072d230a1eb5702e134c971a1fb5724a1c4adf44b9057e461fbf140b7404b \
  'тестове повідомлення' \
  038a4328e0c79249159080f0444edd04a2b22bbdc2c9dbdf6465e201b9eac7faf22f6ab7c4bf94862d9c00851cc6f62da088bfcbba4265fb5297485db0ff9f01
\end{minted}
\begin{minted}{output}
INVALID SIGNATURE!
error: SignatureVerificationFailed
\end{minted}

Коректний підпис для певного повідомлення може бути
сформований лише власником секретного ключа.
Для перевірки підпису застосовується публічний ключ.
Лише для публічного ключа, який відповідає секретному ключу
підписувача, перевірка коректного підпису для коректних
даних буде успішною. Модифікація повідомлення або підпису
виявляються перевіркою.

Таким чином, цифровий підпис забезпечує перевірку
цілісності інформації, а також автентичність відправника.

\chapter{Практичне дослідження узгодження спільного секрету}

Розглянемо узгодження спільного секрету на основі функції
X25519, що реалізує схему Диффі-Геллмана на еліптичній
кривій Curve25519 \cite{rfc7748}.
Програма на Zig v0.17.1, наведена у \ref{listing:3.1},
використовує реалізацію X25519 зі стандартної бібліотеки.

\namedfilelisting{zig}{3.1.zig}
{Програма на Zig для формування спільного секрету за X25519}
\label{listing:3.1}

Формуємо пару ключів \enquote{Alice}:
\begin{minted}{shell}
./3.1.zig generate
\end{minted}
\begin{minted}{output}
Public key: 13fec58c4a50118ec38cedb3cbe8b42930a6b47cfbf056b54526f651c3bfa54c
Private key: dabd6acd9b35c8baa53526d7d857de55b243d33b85584a7dc889c68b08525aec
\end{minted}

Формуємо пару ключів \enquote{Bob}:
\begin{minted}{shell}
./3.1.zig generate
\end{minted}
\begin{minted}{output}
Public key: 4941460421a969164e633bf99554b019a2b35ef8ef7994013eccb00b2944d034
Private key: 86b3944d1da18a37fe1fb2ea697b1516f9046f67b8c6d4c925f512b044600320
\end{minted}

Моделюємо отримання спільного секрета \enquote{Alice},
використовуючи приватний ключ \enquote{Bob}:
\begin{minted}{shell}
./3.1.zig compute  \
  4941460421a969164e633bf99554b019a2b35ef8ef7994013eccb00b2944d034 \
  dabd6acd9b35c8baa53526d7d857de55b243d33b85584a7dc889c68b08525aec
\end{minted}
\begin{minted}{output}
Shared secret: f288d8746c4bc61eb8fc1f2faaff47017193c29a8990cff16fda02357109f922
\end{minted}

Моделюємо отримання спільного секрета \enquote{Bob},
використовуючи приватний ключ \enquote{Alice}:
\begin{minted}{shell}
./3.1.zig compute \
  13fec58c4a50118ec38cedb3cbe8b42930a6b47cfbf056b54526f651c3bfa54c \
  86b3944d1da18a37fe1fb2ea697b1516f9046f67b8c6d4c925f512b044600320
\end{minted}
\begin{minted}{output}
Shared secret: f288d8746c4bc61eb8fc1f2faaff47017193c29a8990cff16fda02357109f922
\end{minted}

Сформований секрет для обох учасників є однаковим,
оскільки в результаті скалярного множення тих самих точок
еліптичної кривої обидва отримують те саме кінцеве
значення. Скалярне множення на еліптичній кривій є
багаторазовим додаванням, а операції, задані на множині
точок, утворюють адитивну групу групу.

Сторони обмінюються лише публічними ключами - значеннями,
за якими обчислювально складно знайти секретний ключ.
Складність ґрунтується на задачі знаходження
дискретного логарифма на еліптичній кривій.
Ані спільний секрет, ані секретні ключі ніколи не
передаються іншим учасникам.

У власне механізмі узгодження спільного секрету немає
засобу підтвердження осои іншої сторони. Якщо сторони
попередньо на обмінялися публічними ключами через
\emph{захищений} канал, вони на матимуть можливості надійно
автентифікувати одна одну під чвс узгодження секрету.
Якщо одна сторона надсилає іншій свій публічний ключ, то
нападник із можливістю підміни повідомлень просто замінить
ключ на свій; тоді інша сторона узгодить спільний секрет з нападником.

\chapter{Вибір криптографічного механізму}

\section{Ситуація 1}

Необхідно використати цифровий підпис, оскільки, зокрема,
необхідне підтвердження джерела підпису.
ECDSA є гарним вибором через забезпечення достатнього рівня захисту
при менших довжинах ключа, оптимальній швидкості обчислення.

\section{Ситуація 2}

Необхідно використати алгоритм узгодження спільного секрету.
Для забезпечення підтвердження особи сторін необхідно
також залучити протокол з використанням статичних ключів.
наприклад, у \cite{rfc8418} пропонується використання
функцій X25519, X448 у схемах узгодження спільного секрету та
для автентифікації відправника.

\section{Ситуація 3}

Через особливості шифрування за асиметричними схемами,
наведеними у частині \ref{chap:asymcrypt-principles},
для шифрування великого обсягу даних краще використати
гібридну криптосистему.

В якості схеми захисту в такій ситуації можна використати
CMS EnvelopedData з використанням ключів на еліптичних
кривих \cite{rfc5753}
якщо не потрібна автентифікація відправника та гарантія
цілісності для відправника, CMS AuthEnvelopedData \cite{rfc5753}
якщо необхідна гарантія цілісності.

\section{Ситуація 4}

Асиметрична криптографія потрібна для створення захищеного, довіреного
каналу між користувачем та сервером для обміну даними під
час автентифікації для запобігання перехоплення паролів чи їх гешів.

Для зберігання можна застосовувати \enquote{традиційне}
гешування із сіллю, або, за можливості -- краще
використовувати спеціалізовані Password-Based Key
Derivation Functions (також із сіллю).

Звісно, для захисту власне записів гешованих паролів можуть
бути застосовані різні схеми, зокрема й із використанням
асиметричної криптографії. Але це залежить від потреб
сервісу в конкретному випадку.

\section{Ситуація 5}

Для розв'язання цієї задачі достатньо використати
механізм цифрового підпису. Публічний ключ виробника
вбудований в поточну версію програмного забезпечення
пристрою. Кожне оновлення підписане виробником за допомогою
безпечно збереженого секретного ключа, і пристрій
перевіряє підпис для відповідного публічного ключа.

\chapter{Аналіз помилок застосування асиметричної криптографії}

\section{Ситуація 1}

Є ризик атак з обраним шифротекстом, що може призвести до
компрометації криптосистеми і розкриття конфіденційних даних,
як показано в \cite{Bleichenbacher1998,cryptoeprint:2017/1189}.

Необхідно використовувати сучасні схеми padding, такі як
RSA-OAEP \cite{rfc8017}.

\section{Ситуація 2}

Повторне використання значення nonce для підпису різних
повідомлень на тій самій ключовій парі призводить до
розкриття секретного ключа, як показано в
\cite{cryptoeprint:2025/654}.

Для запобігання цьому у \cite{fips186-5} наведено безпечні
способи генерації nonce випадковим та детермінованим чином.

\section{Ситуація 3}

При використанні передбачуваного генератора випадкових чисел
можливе виникнення умов як для безпосереднього розкриття
секретного ключа, так і для його розкриття через іншу
доступну інформацію, наприклад, коли в схемі RSA відомі два
ключи, модулі яких не є взаємно простими.
Згідно з \cite{nist-sp-800-133r2}, рівень стійкості захисту
алгоритмів з використанням ключів обмежується також рівнем захисту
процесу генерації ключів. Це, зокрема, встановлює
мінімальний необхідний рівень ентропії, доступний
генератору випадкових бітів.

\section{Ситуація 4}

Довіра до ключа без перевірки його автентичності створює
умови для атак типу \enquote{main-in-the-middle}. Нападник
може модифікувати повідомлення між учасниками взаємодії, та
замінити публічний ключ на свій.

Для протидії цьому ризику
сторони повинні провести обмін публічними ключами через
захищений канал зв'язку заздалегідь, або покладатися на
повноцінну інфраструктуру відкритих ключів, в якій ця
довіра встановлена через третю сторону -- центр сертифікації.

\section{Ситуація 5}

Використання однієї ключової пари у кількох різних
алгоритмах, протоколах, механізмах створює ризик компрометації
безпеки залучених механізмів.
Можливість зниження безпеки через використання ключа у
кількох механізмах,
а також конфлікт між об'єктивними вимогами та кращими
практиками керування ключами в різних процесах зазначені у
\cite{nist-sp-800-57pt1r5}.
Реалізація цих ризиків продемонстрована на практиці
для IPSec IKE у \cite{usenix-27-sec-18-key-reuse-ipsec-ike}
та для застосування RSA стандартом платіжних карт EMV у
\cite{Degabriele2012}.

Отже, безпека криптографічних систем ґрунтується не тільки
на математичній стійкості залучених механізмів, а й
коректному застосуванні і поєднанні цих механізмів.

\chapter{Сучасна постквантова асиметрична криптографія}

\enquote{Класична} асиметрична криптографія, представлена
алгоритмами, що ґрунтуються на складності задач факторизації
цілих чисел та пошуку дискретного логарифму, в тому числі
на еліптичних кривих, є вразливими до квантового алгоритму
факторизації Шора (за реалізації можливості створення
достатньо потужного квантового комп'ютера) \cite{nist-ir-8547}.

Питання захисту від, на даний момент більш теоретичної, ніж
практичної, майбутньої загрози сприймається світовою
спільнотою достатньо серйозно. Відповідні рекомендацій та
нормативні вимоги для підготовки і забезпечення переходу нв
\enquote{пост-квантову}, стійку до атак з використанням
достатньо потужних квантових комп'ютерів,
криптографію, були опубліковані у США \cite{us-eo-14412},
ЄС \cite{eu-ec-cr-pqc}.

Пост-квантова криптографія (Post-Quantum Cryptography,
PQC) включає механізми асиметричної криптографії, які
ґрунтуються на проблемах, що є обчислювально складними
як для звичайних комп'ютерів, так і для квантових, або
принаймні для яких наразі відсутні квантові та інші
ефективні алгоритми розв'язання.

Серед наразі стандартизованих NIST пост-квантових механізмів
наявні Module-Lattice-Based
Key-Encapsulation Mechanism (ML-KEM) -- механізм для
безпечної передачі секретного ключа незахищеним каналом
\cite{nist-fips-203};
Module-Lattice-Based Digital Signature (ML-DSA) -- механізм
цифрового підпису \cite{nist-fips-204}.
Обидва ґрунтуються на складності Задача \enquote{навчання з
помилками} над модулями.

Механізм встановлення ключів (KEM) дозволя учасникам
взаємодій через незахищений канал
обчислити однакове секретне значення таким чином, що воно
буде відомо лише їм \cite{nist-fips-203}.

Цифровий підпис дозволяє
виявити неавторизовані зміни в даних,
автентифікувати підписувача як володаря відповідного секретного ключа.
Цифровий підпис може також слугувати доказом для третьої
сторони того, що підпис був дійсно створений заявленим
підписувачем \cite{nist-fips-204}.

Гібридні механізми, зокрема, гібридне встановлення ключів,
поєднують алгоритми класичної та пост-квантової
криптографії в одному протоколі. Це підсилює захист,
оскільки пост-квантові алгоритми об'єктивно нові й менш
досліджені, ніж класичні. У \cite{nist-ir-8547} та
пов'язаних публікаціях
описан метод поєднання двох або більше механізмів
узгодження ключів, за якого системи залишається безпечною
поки хоча б один з складових механізмів залишається безпечним.

Поступовий перехід на пост-квантові алгоритми зумовлений
тим, що з практичних міркувань здійснення такого переходу
не може бути швидким. Хоча наразі відсутні \enquote{доречні
для криптоаналітичних застосувань квантові комп'ютери},
існує теорема Моска, що
стверджує про необхідність здійснення переходу вже зараз
для захисту від розкриття даних, які зберігаються на
поточний момент часу,
у разі створення квантового комп'ютера, який можливо
застосувати для здійснення атак. Описана в теоремі загроза
отримала назву \enquote{harvest now, decrypt later} (HNDL)
\cite[][\pagename 8]{nist-ir-8547}.
Це особливим актуально щодо інформації державної важливості,
строки захисту якої можуть становити десятки років.

Досі залишаються значні невизначеності у оцінці рівнів
захисту пост-квантових криптосистем, пов'язані із
неможливістю точно оцінити показники майбутніх квантових
комп'ютерів, такі як їх вартість та обчислювальні ресурси
\cite{nist-ir-8547}.

\chapter{Індивідуальне завдання: розділення секрету}

\section{Задача розділення секрету}

Розділення секрету дозволяє розділити ключ або інше
секретне значення у між кількома учасниками так,
що для підновлення ключа необхідно поєднання не менше
заданого числа частин. Такий підхід, зокрема, дозволяє
підсилити  таємність ключа під час зберігання,
оскільки розкриття однією чи небагатьох частин не розкриває
весь секретний ключ.
При цьому в \enquote{пороговій криптографії} (threshold cryptography)
для застосування секретного ключа немає потреби у його
одночасному поєднанні.
Сторони можуть незалежно виконувати обчислення над вхідними даними із
власними чистинами секрету, не розкриваючи їх іншим.
\cite{nist-ir-8214}

\section{Основна ідея розділення секрету}

В описаних у \cite{nist-ir-8214} схемах, розділення секрету
досягається шляхом його представлення як функцію від
кількох аргументів,
де часткове знання аргументів не дає достатньо
інформації для однозначного визначення секрету, а його
підбір за відомими обмеженнями не є обчислювальне
ефективним. В n з n схемах
необхідні всі частини, тоді як в k з n схемах секрет може
бути відновлений за внеском лише частки учасників.
Наприклад, в схемах Шаміра секретом є точка,
задана на поліномі, а частинами є відомі пари аргументів та
інших точок на цьому поліномі. k учасників можуть
однозначно відновити поліном k-порядку, а загалом може бути
згенеровано n частин. Схожа побудова в схемах Блеклі оперує
точками в багатовимірних просторах та перетином гіперплощин.
Для RSA підпису наводиться варіант n з n схеми, де секрет
не розкривається його учасникам, а вони поетапно формують
цифровий підпис. Можливість цього зумовлена математичними
властивостями операцій, які обчислюють значення підпису в
RSA криптосистемі. Розрізняють також схеми із арбітром
(dealer), який знає секрет, та без нього.

\section{Криптографічні примітиви в основі розділення секрету}

Розділення секрету спирається на застосування модульної арифметики;
арифметики у скінченних полях;
генерації випадкових бітів;
а також обчислень, специфічних для конкретної схеми,
доказів з нульовим розголошенням,

Задача розподілення секрету саме по собі вимагає створення
низки криптографічних примітивів для застосування у більш
високорівневих механізмах з властивостями порогової
криптографії \cite{nist-ir-8214}.

\section{Властивості безпеки}

Розділення секрету усуває єдину точку відмови та розподіляє довіру.

З точки зору властивостей інформації цей механізм
забезпечує для секрету захист конфіденційності (зокрема в
схемах, що гарантують теоретико-інформаційну безпеку),
доступності (в k з n схемах). В деяких схемах, як згадане у
\cite[][\pagename iv]{nist-ir-8214} -- цілісності даних
навіть пр компрометації частини секрету.

\section{Практичне застосування}

Практичні застосування включають створення
алгоритмів цифрового підпису, де жоден компонент не володіє
ключем незалежно;
алгоритмів шифрування, які дозволяють зберегти захист даних
навіть якщо один з компонентів скомпрометований та
намагається спотворити результат операцій;
алгоритмів генерації випадкових значень, як працюють
коректно навіть коли деякі донори випадковості ненадійні
або недоступні \cite{nist-ir-8214}.

Серед відомих програмних продуктів, які покладаються на
розділення секрету в пороговій криптографії, можна відзначити:
\begin{enumerate}
  \item
    \href{https://developer.hashicorp.com/vault/docs/concepts/seal}
    {HashCorp Vault}
    -- система управління секретами.
    Кореневий ключ шифрування сховища за схемою Шаміра
    розбивається на кілька частин.
    Для розблокування сховища після перезапуску необхідна
    участь кворуму адміністраторів.
  \item Апаратні ключи, зокрема
    \href{https://trezor.io/trezor-safe-7}{Trezor}
    -- розділення секретної фрази та її резервне зберігання
    на базі розділення секретів за такими специфікаціями,
    як
    \href{https://slips.readthedocs.io/en/latest/slip-0039/}{SLIP-0039}.
\end{enumerate}

\section{Обмеження та ризики при застосуванні розділення секрету}

Наразі для схем розділення секрету актуальними ризиками і
обмеженнями, між якими доводиться шукати компроміс в
практичних застосуваннях, є:

\begin{enumerate}
  \item компроміси між значеннями k, n, та властивостями
    математичних об'єктів в основі схеми;
  \item наявність арбітра (dealer), який знає весь секрет,
    і довіри до нього;
  \item компроміс між безпекою та доступністю.
\end{enumerate}

Для реалізацій механізму характерні джерела ризиків, властиві й
іншим криптографічним засобам:
витоки інформації сторонніи каналами,
якість ентропії при застосування генераторів випадкових бітів
\dots

\chapter{Висновки}

Сучасні асиметричні криптосистеми ґрунтуються
на обчислювальній складності односторонніх функції. Серед
основних проблем в основі класичної асиметричної криптографії
є складність факторизації цілих чисел, складність пошуку
дискретного логарифму. Більш сучасні схеми на основі
еліптичних кривих дозволяють отримати порівняний рівень
захисту при значно менших довжинах ключів та меншій
складності прямого обчислення.

Актуальною загрозою для класичної асиметричної криптографії
є створення квантового комп'ютера, здатного бути
використаним для реалізації атак з ефективними алгоритмами,
такими як алгоритм
факторизації Шора, на примітиви в основі безпеки криптосистем.
Для протидії цьому ризику відбувається розробка та
поступове введення в експлуатацію
пост-квантових криптографічних алгоритмів. В умовах
невизначеності щодо надійності таких алгоритмів та
можливостей квантових обчислень, застосування гібридних
криптосистем є оптимальним підходом.

Безпечне застосування асиметричних криптографічних
алгоритмів повинно враховувати об'єктивні оцінки рівня захисту,
стан розвитку обчислювальної техніки, криптоаналітичної теорії.
Поєднання асиметричних криптографічних механізмів без
детального дослідження їх взаємного впливу створює ризики безпеки.

Дотримання обґрунтованих і підтверджених рекомендацій
провідних дослідницьких закладів та державних інституцій
зокрема таких як NIST, IETF, ANNSSI,
та поступовий перехід до пост-квантової криптографії з
застосуванням гібридних криптосистем є найкращими практиками
захисту даних за допомогою асиметричної криптографії.

\printbibliography[
  heading=bibintoc,
  title={Перелік використаних джерел},
]

\end{document}
